% !PRINTED 749
\weberpage{749}

% [sic] -- the source uses $\delta_1$ and $\delta_2$ both for scalar lengths and for intervals; the printed symbols are preserved.
We denote by $\alpha_1$, $\beta_1$ the two \gesperrt{earliest} numbers of the sequence $\Omega$ that lie in the interval $\delta$, we assume $\alpha_1 < \beta_1$, and we set $\delta_1 = \beta_1 - \alpha_1$, so that $\delta_1 = (\alpha_1, \beta_1)$ is a part of the interval $\delta$.

Now in the same way we denote by $\alpha_2$, $\beta_2$ the two earliest numbers of $\Omega$ that lie \gesperrt{in the interior} of the interval $\delta_1$ (with the exclusion of the limits), we presuppose $\alpha_2 < \beta_2$, and we set $\delta_2 = \beta_2 - \alpha_2$.

In this way we can continue, and we obtain an unlimited sequence of intervals:
\[
\delta,\ \delta_1,\ \delta_2,\ \delta_3,\ \dots,
\]
each of which includes all the following ones, and two sequences of numbers:
\[
\begin{gathered}
\alpha,\ \alpha_1,\ \alpha_2,\ \dots \\
\beta,\ \beta_1,\ \beta_2,\ \dots
\end{gathered}
\]
which, with the possible exception of the first, $\alpha$ and $\beta$, all belong to the sequence $\Omega$. The $\alpha, \alpha_1, \alpha_2, \dots$ form an increasing, the $\beta, \beta_1, \beta_2, \dots$ a decreasing sequence of numbers, and at the same time every $\alpha$ is smaller than every $\beta$.

1. \gesperrt{From this it follows that the numbers $\alpha_\nu$ have an upper limit $a$, and the numbers $\beta_\nu$ a lower limit $b$, and that $a$ is in any case not greater than $b$.} (Cf. Bd. I, \S.\,38, 1.)

But it is possible that $a = b$.

From the mode of formation of the intervals $\delta, \delta_1, \delta_2, \dots$ the following further results:

2. \gesperrt{If any number $\omega$ of the sequence $\Omega$ lies in the interval $\delta_\nu$, with the exclusion of its limits $\alpha_\nu, \beta_\nu$, then $\omega$ stands later in the sequence $\Omega$ than the pair of numbers $\alpha_\nu, \beta_\nu$ belonging to that same sequence.}

For $\alpha_\nu, \beta_\nu$ were indeed the two \gesperrt{earliest} numbers of $\Omega$ lying in the interval $\delta_{\nu-1}$. From this it follows:

3. \gesperrt{The pairs of numbers $\alpha_\nu, \beta_\nu$ belonging to the sequence $\Omega$ are the later terms of the sequence $\Omega$ the greater the index $\nu$ is, and since we have assumed that the sequence of intervals $\delta_\nu$ does not break off, we can, for a sufficiently large $\nu$, let $\alpha_\nu, \beta_\nu$ move out arbitrarily far in the sequence $\Omega$.}
